% !TEX root = ../main.tex
\chapter{Background}
\label{ch:background}

\section*{Outline of the chapter}

This chapter collects the technical ingredients on which the thesis
builds: regular variation and the second-order condition; the Hill and
Weissman estimators for the tail index and extreme quantiles; expectiles
and their characterisation as a coherent and elicitable risk measure;
single-sample iid extreme-expectile inference; and the optimal pooling
framework of \textcite{DaouiaPadoanStupfler2021}. Nothing in this
chapter is original; the contributions begin in \Cref{ch:pool-int}.

\section{Regular variation and the second-order condition}
\label{sec:bg:rv}

Throughout the thesis we work with a real-valued random loss $X$ that
has continuous distribution function $F$ and survival function
$\overline F = 1 - F$. Given an iid sample $X_{1}, \ldots, X_{n}$ from
$F$, the corresponding order statistics are denoted by
$X_{1:n} \le X_{2:n} \le \cdots \le X_{n:n}$.

Heavy-tailed asymptotics are governed by the language of regular
variation; the reference is \textcite[Ch.~1]{deHaanFerreira2006}.

\begin{definition}[Regular variation]
\label{def:bg:rv}
A measurable function $f : \R_{+} \to \R_{+}$ is \emph{regularly varying
at infinity with index} $\alpha \in \R$, written
$f \in \mathrm{RV}_{\alpha}$, if for every $y > 0$,
\[
  \lim_{t \to \infty} \frac{f(t y)}{f(t)} \;=\; y^{\alpha}.
\]
The functions in $\mathrm{RV}_{0}$ are called \emph{slowly varying}.
\end{definition}

The first-order condition under which all of our asymptotics are taken
asserts that $F$ has a Pareto-type tail.

\begin{condition}[First-order condition]
\label{cond:bg:first-order}
There exists $\gamma > 0$ such that
$\overline F \in \mathrm{RV}_{-1/\gamma}$, or equivalently the
\emph{tail quantile function}
\[
  U(t) \;\coloneqq\; \inf\!\Big\{ x \in \R \,:\, \tfrac{1}{\overline F(x)} \ge t \Big\}
\]
satisfies $U \in \mathrm{RV}_{\gamma}$.
\end{condition}

\Cref{cond:bg:first-order} is equivalent to saying that $X$ lies in the
max-domain of attraction of the Fr\'echet distribution with shape
$\gamma$ \parencite[Theorem~1.2.1]{deHaanFerreira2006}. The parameter
$\gamma > 0$ is the \emph{extreme-value index} of $F$; it controls how
slowly $\overline F$ decays at infinity. Standard heavy-tailed
distributions encountered in risk management satisfy
\Cref{cond:bg:first-order}, including the Pareto, Burr, Student-$t$,
Fr\'echet, and $F$-distribution families
\parencite[\S2.6]{deHaanFerreira2006}.

For the asymptotic-normality statements that drive the inferential
theory we need a quantitative refinement of \Cref{cond:bg:first-order}
controlling the rate of convergence in
\Cref{def:bg:rv}. The standard refinement is the second-order
condition.

\begin{condition}[Second-order condition $\mathcal C_{2}(\gamma, \rho, A)$]
\label{cond:bg:second-order}
There exist $\gamma > 0$, a second-order parameter $\rho \le 0$, and a
measurable auxiliary function $A$ of constant sign with $A(t) \to 0$ as
$t \to \infty$, such that for every $y > 0$,
\[
  \lim_{t \to \infty} \frac{1}{A(t)}
    \!\left[ \frac{U(t y)}{U(t)} - y^{\gamma}\right]
  \;=\; y^{\gamma}\, \frac{y^{\rho} - 1}{\rho},
\]
where the right-hand side is read as $y^{\gamma} \log y$ when
$\rho = 0$.
\end{condition}

The function $|A|$ is itself regularly varying with index $\rho$
\parencite[Theorem~2.3.9]{deHaanFerreira2006}; smaller $|\rho|$ means
slower convergence and therefore larger second-order bias of the
tail-index estimators we shall encounter in
\Cref{sec:bg:tail-index-quantile}.

For the expectile--quantile asymptotic equivalence
(\Cref{sec:bg:extreme-expectile}) we will impose the stronger range
$\gamma \in (0, 1/2)$, which guarantees both that $X$ has finite second
moment and that the expectile $\expectile{\tau}(X)$ admits a
non-degenerate $\tau \uparrow 1$ limit when scaled by the corresponding
quantile $\Quant_{X}(\tau)$.

\section{Tail-index and extreme-quantile estimation}
\label{sec:bg:tail-index-quantile}

Under the first-order condition, the extreme-value index $\gamma > 0$
is the primary parameter determining the behaviour of $F$ in the tail.
We collect here the two classical estimators that the thesis builds on:
the Hill estimator of $\gamma$ and the Weissman extrapolation for
extreme quantiles. Both are standard; the textbook reference is
\textcite[Ch.~3--4]{deHaanFerreira2006}.

\subsection{The Hill estimator}
\label{sec:bg:hill}

Given an iid sample $X_{1}, \ldots, X_{n}$ from $F$ and an integer
$k = k_{n} \in \{1, \ldots, n-1\}$, the \emph{Hill estimator} of $\gamma$
based on the $k+1$ top order statistics is
\begin{equation}
\label{eq:bg:hill}
  \hat\gamma_{n}(k) \;=\; \frac{1}{k} \sum_{i=1}^{k} \log\!\left(
    \frac{X_{n-i+1\,:\,n}}{X_{n-k\,:\,n}} \right).
\end{equation}
The interpretation is direct: if the upper tail of $F$ were exactly
Pareto with index $1/\gamma$, the log-spacings
$\log(X_{n-i+1:n}) - \log(X_{n-k:n})$ would be iid exponential with
mean $\gamma$, and $\hat\gamma_{n}(k)$ would be their sample mean.
\Cref{cond:bg:first-order} asserts that the upper tail is
\emph{approximately} Pareto, which translates into the following
asymptotic theory \parencite[Theorem~3.2.5]{deHaanFerreira2006}.

\begin{theorem}
\label{thm:bg:hill}
Assume $\mathcal C_{2}(\gamma, \rho, A)$ with $\gamma > 0$. Let
$k = k_{n} \to \infty$ and $k / n \to 0$ as $n \to \infty$. Then
$\hat\gamma_{n}(k) \toprob \gamma$. If, in addition,
$\sqrt{k}\, A(n/k) \to \lambda \in \R$, then
\[
  \sqrt{k}\,\big( \hat\gamma_{n}(k) - \gamma \big)
  \;\todist\;
  \mathcal N\!\left( \frac{\lambda}{1 - \rho},\; \gamma^{2} \right).
\]
\end{theorem}

The choice of $k$ controls the well-known bias--variance trade-off: too
small a $k$ inflates the variance, while too large a $k$ injects
second-order bias through the term $\lambda / (1-\rho)$. Adaptive choice
of $k$ is a research topic in its own right; we do not pursue it in
this thesis and treat $k$ as fixed by the data analyst.

\subsection{Weissman extrapolation to extreme quantiles}
\label{sec:bg:weissman}

The \emph{Weissman extrapolation} \parencite{Weissman1978} extrapolates
from the intermediate-order statistic $X_{n-k:n}$ --- which estimates
the quantile $\Quant_{X}(1 - k/n)$ at the intermediate level $1 - k/n$
--- to a quantile at a very-extreme level $1 - p$ with $p = p_{n} \to 0$
as $n \to \infty$. Plugging the Hill estimator into the first-order
identity $U(t y) / U(t) \to y^{\gamma}$ gives the Weissman estimator
\begin{equation}
\label{eq:bg:weissman}
  \hatWei_{n}(1 - p \mid k)
  \;=\; \left( \frac{k}{n p} \right)^{\hat\gamma_{n}(k)}
        X_{n - k\,:\,n}.
\end{equation}
The typical regime of interest is $n p$ bounded, i.e.\ extrapolation
beyond the largest observation in the sample.

\begin{theorem}
\label{thm:bg:weissman}
Assume $\mathcal C_{2}(\gamma, \rho, A)$ with $\gamma > 0$ and
$\rho < 0$. Let $k = k_{n} \to \infty$, $k/n \to 0$, $p = p_{n} \to 0$
with $k/(np) \to \infty$ and $\sqrt{k}/\log(k/(np)) \to \infty$.
Suppose $\sqrt{k}\, A(n/k) \to \lambda \in \R$. Then
\[
  \frac{\sqrt{k}}{\log\!\big(k/(np)\big)}
  \left( \frac{\hatWei_{n}(1 - p \mid k)}{\Quant_{X}(1 - p)} - 1 \right)
  \;\todist\;
  \mathcal N\!\left( \frac{\lambda}{1 - \rho},\; \gamma^{2} \right).
\]
\end{theorem}

The rate $\sqrt{k}/\log(k/(np))$ reflects the cost of extrapolation:
the further into the tail one extrapolates, the larger the
log-denominator and the slower the rate. The asymptotic variance is
identical to that of the Hill estimator (\Cref{thm:bg:hill}); the
extrapolation contributes only through the multiplicative log-scale
factor. This is the structural fact that
\textcite[\S2.3]{DaouiaPadoanStupfler2021} exploits to motivate the
weighted-\emph{geometric} pooling of Weissman estimators --- rather than
the arithmetic pooling natural for centred quantities --- and that we
shall in turn exploit in \Cref{ch:pool-ext} for the extrapolation of
pooled expectiles.

\section{Expectiles}
\label{sec:bg:expectiles}

\todo{Definition via the asymmetric-LS minimisation of
\textcite{NeweyPowell1987}; first-order conditions for the population
expectile; basic properties (existence for $\E|X|<\infty$, uniqueness,
location/scale equivariance, monotonicity in $\tau$, $\xi_{1/2}=\E X$).
Coherence (for $\tau \ge 1/2$) and elicitability via the asymmetric
squared scoring function; cite
\textcite{BelliniKlarMullerRosazzaGianin2014}. Choquet representation
of \textcite{EhmGneitingJordanKrueger2016}. Risk-management
interpretation as a VaR-distortion.}

\section{Extreme expectile inference: the iid heavy-tailed baseline}
\label{sec:bg:extreme-expectile}

\todo{State and reference the expectile--quantile asymptotic
equivalence: $\xi_\tau/Q_X(\tau) \to (\gamma^{-1}-1)^{-\gamma}$ as
$\tau\uparrow 1$ for $\gamma\in(0,1)$
\parencite[Proposition~1]{DaouiaGirardStupfler2019}.}

\todo{The two estimator families in the iid setting
\parencite{DaouiaGirardStupfler2018}: (a) LAWS direct,
$\tildexpectile{\tau_n} = \argmin_{\theta}\sum_i \eta_{\tau_n}(X_i-\theta)$,
with asymptotic variance $2\gamma^3/(1-2\gamma)$ at intermediate levels;
(b) QB plug-in, $\hatexpectile{\tau_n} = (\hat\gamma^{-1}-1)^{-\hat\gamma}
\hatQuant(\tau_n)$, with asymptotic variance involving
$m(\gamma)=(1-\gamma)^{-1}-\log(\gamma^{-1}-1)$.}

\todo{Weissman-style extrapolation to very-extreme levels $\tau_n^\prime$
with $n(1-\tau_n^\prime) \to c$: $\overline{\xi}^\star_{\tau_n^\prime} =
((1-\tau_n^\prime)/(1-\tau_n))^{-\hat\gamma}\,\overline{\xi}_{\tau_n}$,
with asymptotic variance $\gamma^2$ in the iid case. The LAWS-anchored
and QB-anchored extrapolations have the same asymptotic distribution at
the extreme level (the tail-index estimator dominates).}

\todo{State the plug-in CLT of
\textcite[Proposition~C.6]{DavisonPadoanStupfler2023}, which is the
abstract technical lemma that lifts joint $\sqrt{n(1-\tau_n)}$-Gaussian
limits of $(\hat\gamma, \overline\xi)$ to limits of the extrapolated
$\overline\xi^\star$. This is the workhorse for \Cref{ch:pool-int} and
\Cref{ch:pool-ext}.}

\section{The optimal pooling framework}
\label{sec:bg:pooling}

\todo{Setup: $m$ samples $\bm{X}_1,\ldots,\bm{X}_m$, sample sizes $n_j$,
total $n = \sum n_j$; effective sample sizes $k_j$ for tail estimation;
ratios $n_j/n \to b_j$, $k_j/n_j \to 0$, $k_j/k \to c_j$. Marginal
second-order condition $\mathcal C_2(\gamma_j,\rho_j,A_j)$. Tail-copula
assumption $\mathcal J(\bm R)$ encoding cross-sample asymptotic
dependence via $R_{j,\ell}$. The common-tail case
$\bm\gamma = \gamma\bm 1$ that the thesis adopts.}

\todo{Pooled Hill estimator $\hat\gamma_n(\bm\omega) =
\bm\omega^\top\hat{\bm\gamma}_n$ with $\bm\omega^\top\bm 1 = 1$. State
Theorem~1 of \textcite{DaouiaPadoanStupfler2021}: joint
$\sqrt{k_j}(\hat\gamma_j(k_j)-\gamma_j) \todist \mathcal N(\bm B,\bm V)$,
with bias vector $\bm B$ involving $\lambda_j/(1-\rho_j)$ and covariance
matrix $\bm V$ whose diagonal entries are $\gamma_j^2$ and off-diagonal
entries are
$\gamma_j\gamma_\ell\sqrt{c_\ell/c_j}\,R_{j,\ell}(c_\ell/c_j,b_\ell/b_j)$.}

\todo{Variance-optimal and AMSE-optimal weights:
$\bm\omega^{\mathrm{Var}} = \Vc^{-1}\bm 1/(\bm 1^\top\Vc^{-1}\bm 1)$, and
the AMSE-optimal counterpart involving $\Bc$ and $\Vc$.}

\todo{Weighted geometric pooled Weissman estimator
$\hatWei_n(1-p\mid \bm\omega) = \prod_j [\hat q^\star_j(1-p\mid k_j)]^{\omega_j}$
and Theorem~2 of \textcite{DaouiaPadoanStupfler2021} giving its
asymptotic distribution.}

\todo{Distributed-inference specialisation: iid data within and across
machines forces $R_{j,\ell}\equiv 0$, so $\bm V$ becomes diagonal with
entries $\gamma^2/c_j$; the variance-optimal weights reduce to
$\omega^{\mathrm{Var}}_j \propto k_j$ (effective-sample-size weighting);
and the pooled Hill and Weissman estimators are asymptotically
equivalent to the unfeasible benchmark Hill and Weissman estimators
built from the combined sample.}

\todo{Tail-homogeneity test $H_0:\bm\gamma=\gamma\bm 1$ via the
chi-square statistic $\Lambda_n =
k(\hat{\bm\gamma}_n-\bar\mu_n\bm 1)^\top\hatVc^{-1}
(\hat{\bm\gamma}_n-\bar\mu_n\bm 1)$. Asymptotic confidence interval for
$\gamma$ under tail homogeneity.}
